% Generated draft: requires mathematical and source review.
\documentclass[11pt]{article}
\usepackage[margin=25mm]{geometry}
\usepackage{fontspec}
\setmainfont{Latin Modern Roman}
\usepackage{amsmath,amssymb,mathtools,unicode-math}
\setmathfont{Latin Modern Math}
\usepackage{enumitem,needspace,xurl,hyperref}
\setlength{\parindent}{0pt}
\setlength{\parskip}{5pt}
\setlength{\emergencystretch}{3em}
\newcommand{\N}{\mathbb{N}}
\newcommand{\Z}{\mathbb{Z}}
\newcommand{\Q}{\mathbb{Q}}
\newcommand{\R}{\mathbb{R}}
\newcommand{\C}{\mathbb{C}}
\newcommand{\F}{\mathbb{F}}
\newcommand{\A}{\mathbb{A}}
\newcommand{\PP}{\mathbb{P}}
\newcommand{\E}{\mathbb{E}}
\newcommand{\eps}{\varepsilon}
\newcommand{\dd}{\,\mathrm{d}}
\begin{document}
\section*{UnsolvedMath ID: 701776}
\textbf{Source problem code:} catalogue entry
\subsection{Definitions and mathematical statement}
Let \(p\) be a prime. A pro-\(p\) group is a topological group isomorphic to an inverse limit of finite \(p\)-groups. Write \(H^i_{\mathrm{cts}}(G,\F_p)\) for continuous cohomology with trivial discrete coefficients, and let \(\operatorname{cd}_p(G)\) denote its continuous \(p\)-cohomological dimension. In the usual pro-\(p\) sense, a Poincar\'e duality group of dimension \(3\) has \(\operatorname{cd}_p(G)=3\), finite-dimensional \(H^i_{\mathrm{cts}}(G,\F_p)\) for every \(i\), and \(H^3_{\mathrm{cts}}(G,\F_p)\cong\F_p\), with perfect cup-product pairings
\[
H^i_{\mathrm{cts}}(G,\F_p)\times H^{3-i}_{\mathrm{cts}}(G,\F_p)
\longrightarrow H^3_{\mathrm{cts}}(G,\F_p),\qquad 0\leq i\leq 3.
\]
Let \(F_2\) denote the free pro-\(p\) group on two generators, and give \(F_2\times F_2\) its product topology.

\textbf{Question (P. Zalesskii).} Can a pro-\(p\) Poincar\'e duality group \(G\) of dimension \(3\) contain a subgroup topologically isomorphic to \(F_2\times F_2\)? Equivalently, can such a \(G\) admit a continuous injective homomorphism
\[
F_2\times F_2\hookrightarrow G?
\]
Such an image is automatically closed. This draft uses the usual topological interpretation of containment in a pro-\(p\) group; the supplied statement does not explicitly specify closedness.

The supplied source further states that an affirmative answer would give a negative answer to Kourovka Notebook Problem 20.117. The statement of 20.117 is not supplied, so this implication is retained as a source qualification without reconstruction. The source also reports that the analogous containment is impossible in the abstract-group setting, citing P. H. Kropholler and M. A. Roller; in that comparison, the factors are discrete free groups of rank \(2\), rather than free pro-\(p\) groups. This reported background result is not independently verified here.
\subsection{Short English statement}
P. Zalesskii asks whether a three-dimensional pro-\(p\) Poincar\'e duality group can contain the direct product of two free pro-\(p\) groups of rank \(2\). Containment is interpreted topologically. The supplied source links an affirmative answer to Problem 20.117 and reports a negative result for the analogous abstract-group question.
\subsection{Sources}
\begin{itemize}
\item Local UnsolvedMath registry, permanent ID \(701776\), exact source code \texttt{catalogue entry}, title \emph{Kourovka Notebook Problem 20.118}. The supplied record labels the problem \emph{open}; this is supplied registry status, not an independently verified literature assessment.
\item \emph{The Kourovka Notebook (Unsolved Problems in Group Theory)}, 20th issue (2022), Problem 20.118, proposed by P. Zalesskii. The supplied background also identifies the 21st issue edition (2026), \texttt{arXiv:1401.0300v48}. The collection description names E. I. Khukhro and V. D. Mazurov as editors.
\item P. H. Kropholler and M. A. Roller, \emph{J. London Math. Soc. (2)} \textbf{39} (1989), 271--284. This reference is supplied for the reported impossibility in the abstract-group case; no article title is provided.
\item Kourovka Notebook Problem 20.117 is referenced by the supplied statement for the claimed consequence of an affirmative answer, but its statement is not included.
\end{itemize}
\end{document}
